\newpage
\section{Обзор предметной области}
\label{ch:review}

Система кластеризации научных статей и построенный поверх неё рекомендатель соавторов стоят на
стыке нескольких областей: библиометрии и анализа научных графов, обработки естественного языка,
тематического моделирования и информационного поиска. Чтобы обосновать инженерные решения главы~2,
ниже разобраны пять блоков: источники данных и графовые методы, текстовые эмбеддинги, тематическое
моделирование, рекомендация научных соавторов и применение больших языковых моделей в роли
оценщика. В каждом блоке отмечено, почему «классический» подход в нашей задаче упирается в потолок,
и какой инструмент приходит ему на смену.

\subsection{Научные графы, источники данных и предсказание связей}

Исторически анализ науки строится на графах. Граф цитирований связывает работы по ссылкам, граф
соавторства — исследователей по совместным публикациям. Структуру таких сетей подробно изучил
Ньюман~\cite{newman2001coauthorship}: у большинства авторов соавторов немного, тогда как у
отдельных — очень много, и при этом почти все исследователи оказываются связаны в одну крупную
компоненту. Такие графы хорошо отвечают на вопрос «кто с кем уже связан», но для рекомендации
новых связей их приходится достраивать.

Сырьём для графов служат библиографические базы. Долгое время эталоном был Microsoft Academic
Graph~\cite{sinha2015mag}, однако после его закрытия открытым преемником стал
OpenAlex~\cite{priem2022openalex} — полностью открытый индекс работ, авторов, площадок и
организаций, который мы и используем как источник корпуса.

Задача рекомендации новой связи в графовой постановке — это \emph{предсказание связей} (link
prediction). Либен-Ноуэлл и Клейнберг~\cite{libennowell2007link} показали, что по одной лишь
топологии сети (общие соседи, мера Адамика–Адара и т.\,п.) будущие рёбра предсказываются заметно
лучше случайного, но точность остаётся невысокой. Главное ограничение для нас в другом: чисто
структурные признаки не видят \emph{смысла} работ. Два исследователя могут заниматься почти
одинаковыми задачами, но если в графе между ними нет коротких путей — топологический метод их не
сблизит. Отсюда вытекает потребность в содержательном, текстовом представлении публикаций.

Современные графовые нейросети ослабляют это ограничение: они агрегируют не только структуру
соседства, но и признаки узлов. В GraphSAGE~\cite{hamilton2017graphsage} представление узла
строится через сообщения от соседей и собственные признаки узла, поэтому метод можно применить к
графу соавторства, инициализировав авторов семантическими признаками по их статьям. Такая модель
не заменяет текстовые эмбеддинги, а проверяет, добавляет ли граф прошлых соавторств полезный сигнал
поверх семантического представления автора.

\subsection{Текстовые эмбеддинги}

Чтобы сравнивать статьи по смыслу, текст нужно превратить в вектор. Ранние подходы — мешок слов и
TF-IDF — игнорируют порядок и синонимию, а первые плотные представления слов
(word2vec~\cite{mikolov2013word2vec}) дали векторы для отдельных слов, но не для предложений и тем
более не учитывали контекст: слово получало один вектор на все смыслы.

Перелом принесли трансформерные модели. BERT~\cite{devlin2019bert} — предобученный двунаправленный
энкодер — научился строить \emph{контекстные} представления токенов, однако его выходы в исходном
виде плохо подходят для прямого сравнения текстов по косинусной близости. Эту проблему решил
Sentence-BERT~\cite{reimers2019sbert}: дообучение на парах предложений даёт эмбеддинги, которые уже
можно осмысленно сравнивать скалярным произведением. С тех пор стандартом стали специализированные
энкодеры предложений и документов.

Для нашего корпуса критична многоязычность — статьи OpenAlex написаны на разных языках. Поэтому
для векторизации статей выбрана модель multilingual-e5-base~\cite{wang2024multilinguale5}: она
строит единое векторное пространство для многих языков при префиксе \texttt{query:} и позволяет
кластеризовать и сравнивать работы независимо от языка оригинала. Косинусная близость таких
эмбеддингов служит рабочей мерой семантического сходства на этапах кластеризации и retrieval.
Для метадокументов метакластеров (короткие текстовые описания тем) используется та же
модель multilingual-e5-base с префиксом \texttt{query:}.

\subsection{Тематическое моделирование}

Группировка статей по темам — классическая задача. Долгое время её решало вероятностное
тематическое моделирование, прежде всего латентное размещение Дирихле
(LDA)~\cite{blei2003lda}: документ моделируется как смесь тем, тема — как распределение над словами.
LDA работает с мешком слов и потому плохо переносит короткие тексты (а аннотация — это короткий
текст) и не использует семантику эмбеддингов.

Современная альтернатива — embedding-based тематическое моделирование, канонический представитель
которого BERTopic~\cite{grootendorst2022bertopic}. Его конвейер: эмбеддинги $\to$ понижение
размерности $\to$ плотностная кластеризация $\to$ интерпретируемое представление тем. На шаге
понижения размерности применяется UMAP~\cite{mcinnes2018umap}, сохраняющий локальную структуру
многообразия лучше линейных методов; кластеризацию выполняет HDBSCAN~\cite{campello2013hdbscan,
mcinnes2017hdbscanjoss} — иерархический плотностный алгоритм, который не требует заранее задавать
число кластеров $K$, помечает разреженные точки как шум и устойчив к неравномерной плотности. Это
важные свойства: число тем в корпусе заранее неизвестно, а «фоновые» статьи честнее пометить шумом,
чем приписывать к ближайшей теме силой.

Представление темы словами строится через \emph{class-based} TF-IDF (c-TF-IDF): кластер
рассматривается как один большой документ, и вес слова показывает его специфичность для темы в
целом. Чтобы итоговый список ключевых слов не состоял из синонимов, применяется Maximal Marginal
Relevance~\cite{carbonell1998mmr}, балансирующий релевантность и новизну отбираемых терминов.

Для честного сравнения с плотностным подходом в главе~3 добавлены наивные бейзлайны: случайное
разбиение и K-Means на том же пространстве UMAP-10 с тем же числом кластеров $K$. Они показывают,
насколько выигрыш HDBSCAN по LLM-когерентности и различимости оправдан относительно геометрически
простого разбиения.

\subsection{Рекомендация научных соавторов}

Рекомендация работ и соавторов — отдельное направление; обзор Бееля и
соавторов~\cite{beel2016survey} систематизирует подходы к рекомендации научных публикаций.
Условно методы делятся на два семейства. \emph{Графовые} (по сути — то же предсказание связей)
опираются на структуру сети соавторства и страдают теми же ограничениями: холодный старт и
слепота к смыслу. \emph{Контентные} сравнивают исследователей по текстам их работ — и именно сюда
ложатся эмбеддинги.

С масштабом корпуса прямое попарное сравнение становится дорогим, поэтому в информационном поиске
закрепилась двухэтапная схема \emph{retrieve-then-rerank}. На первом этапе быстрый плотный поиск
(dense retrieval) по близости векторов отбирает небольшой набор кандидатов; парадигму плотного
поиска для открытых вопросов формализовали Карпухин и соавторы в DPR~\cite{karpukhin2020dpr}. На
втором этапе кандидаты переранжируются более дорогой и точной моделью.

Чтобы получить вектор \emph{автора}, а не отдельной статьи, историю его публикаций нужно
агрегировать. Естественный инструмент — трансформерный энкодер~\cite{vaswani2017attention}: он
обрабатывает последовательность эмбеддингов статей автора и через механизм внимания сводит её в
единое представление. Такой обучаемый агрегатор и образует первый, retrieval-этап нашего
рекомендателя. Второй проверяемый вариант — GraphSAGE поверх графа соавторства: он использует
те же семантические признаки авторов, но дополнительно передаёт информацию по рёбрам прошлых
совместных публикаций.

\subsection{Большие языковые модели в роли оценщика}

Появление больших языковых моделей (LLM), способных к обучению по нескольким примерам и следованию
инструкциям~\cite{brown2020gpt3}, изменило не только генерацию текста, но и способ оценки. Там, где
ручная разметка релевантности слишком дорога, всё чаще применяют подход \emph{LLM-as-a-judge}:
языковая модель оценивает качество ответов по заданной шкале. Чжэн и соавторы~\cite{zheng2023judging}
показали, что такие оценки сильно коррелируют с человеческими, и одновременно описали характерные
смещения (позиционное, в пользу многословия), которые приходится контролировать продуманным
промптом и перемешиванием порядка кандидатов.

В нашей задаче LLM используется в трёх ролях: разметка тем человекочитаемыми названиями, оценка
релевантности пар «автор\,--\,кандидат» на этапе переранжирования и генерация объяснений
рекомендаций. Подход LLM-as-a-judge при этом служит и инструментом \emph{измерения} качества там,
где фактические соавторства дают слишком разреженный и заниженный сигнал (см. главу~3).

\subsection{Выводы и результаты по главе}

Обзор показывает логику дальнейших решений. Графовые методы хорошо описывают уже состоявшиеся
связи, но плохо подсказывают новые и не видят смысла работ, поэтому в основу системы кладутся
многоязычные эмбеддинги multilingual-e5-base. Для интерпретируемой тематической структуры выбран
embedding-based подход в стиле BERTopic с плотностной кластеризацией HDBSCAN; качество
сопоставляется с наивными бейзлайнами K-Means и случайного разбиения. Рекомендация соавторов
строится по схеме retrieve-then-rerank: mean-бейзлайн, обучаемый трансформерный агрегатор и
GraphSAGE-агрегатор графа соавторства извлекают кандидатов, а большая языковая модель выполняет
переранжирование и объяснение. Наконец, подход
LLM-as-a-judge выбран как способ честно измерить качество рекомендаций в условиях разреженного
целевого сигнала. Эти выборы и реализуются в главе~2.
